\documentclass{article}
\usepackage[utf8]{inputenc}
\usepackage[spanish]{babel}
\usepackage{hyperref}
\title{Plataforma de Reportes - FES Aragón}
\author{Equipo de Desarrollo}
\date{\today}
\begin{document}
\maketitle
\section{Resumen}
Aplicación para gestionar reportes de incidencias en la FES Aragón. Usuarios con rol \textbf{admin} y \textbf{alumno}.

\section{Arquitectura}
El proyecto está organizado en carpetas \texttt{public/}, \texttt{src/}, \texttt{templates/} y \texttt{database/}.

\section{Despliegue}
Instrucciones: importar \texttt{database/schema.sql} y \texttt{database/seed.sql}. Levantar servidor PHP/Apache o usar Docker.

\section{API / Rutas}
\begin{itemize}
  \item /public/login.php - Iniciar sesión
  \item /public/register.php - Registro de usuarios
  \item /public/report_list.php - Listado de reportes
  \item /public/report_create.php - Crear/editar reporte
  \item /public/report_view.php?id= - Ver detalles
\end{itemize}

\end{document}
